% !PRINTED 207
\weberpage{207}
%<<CONT
\[
\xi = \lambda\xi', \quad \eta = \frac{1}{\lambda}\eta',
\]
whose determinant is $= 1$, leaves the normal form \S.\,65, (4), and in it
$a_2'$, unchanged, while $a_1'$, $a_3'$ pass into $\lambda^2 a_1'$, $a_3' : \lambda^2$,
and here $\lambda$ is an arbitrary quantity.

If, then, we set
\[
a_2' = z
\]
we have
\weq{4}{
I = \varphi(a_1'a_3', z),
}
where $\varphi$ is an integral, rational function with integer coefficients of the two
arguments $a_1'a_3'$ and $z$.

But now, by \S.\,65, (11),
\[
3a_1'a_3' = z^2 - A,
\]
and if we therefore multiply (4) by a suitable power of 3, say $3^\nu$,
we obtain
\weq{5}{
3^\nu I = \psi(A, z),
}
where $\psi$ is again a function of $A$ and $z$ with integer coefficients. But by
\S.\,65, (12),
\[
z^3 = 3Az - B,
\]
and by means of this we can express all higher powers of $z$ through the first and
second, and thus obtain
\weq{6}{
3^\nu I = \chi(A, B, z),
}
where $\chi$ is at most of the second degree with respect to $z$.

But the left side of (6) remains unchanged when $z$ is replaced by $a_2'$,
$a_2''$, $a_2'''$, that is, by three different values, and
consequently the function $\chi$ cannot contain the variable $z$ at all any longer
(\S.\,30, II).

Therefore
\weq{7}{
3^\nu I = \chi(A, B),
}
where $\chi$ is an integral, rational function with integer coefficients.

If in (7), by (3), we set
\[
B^2 = 4A^3 - 27 D
\]
it follows that (7) has one of the two following forms
\weq{8}{
3^\nu I = \Phi(A, D), \quad B\Phi(A, D),
}
according as the degree of $I$ is even or odd.
